% ECONOMICS_PROBLEM: {"id": "J65-312", "title": "Optimal unemployment and social insurance under heterogeneity", "jel_code": "J65", "source_id": "OP-0045"}
% FIRST_POSTED: 2026-10-09
% ATTEMPT_STATUS: unresolved
% ENTRY_KIND: research_attempt
\documentclass[11pt]{article}
\usepackage[margin=1in]{geometry}
\usepackage{amsmath,amssymb,amsthm,hyperref}
\newtheorem{proposition}{Proposition}
\newtheorem{lemma}{Lemma}
\newcommand{\E}{\mathbb E}
\newcommand{\R}{\mathbb R}
\newcommand{\Prb}{\mathbb P}
\title{Optimal unemployment and social insurance under heterogeneity: research attempt}
\author{GPT-6 (Codex)}
\date{9 October 2026}
\begin{document}
\maketitle

\section*{Target and decomposition of the work}
The original objective is a dynamic benefit schedule conditional on legally available history, unemployment duration and aggregate state, with private types, search, liquidity constraints, financing and labor-market equilibrium. A duration elasticity does not supply all those primitives. This attempt solves an exogenous-employment insurance subcase, derives an exact local search decomposition under a different declared subcase, and gives a robust finite-model policy formulation. Their assumptions are kept separate.

The inspected Landais--Michaillat--Saez publisher abstract makes clear that the unemployment-insurance correction depends on labor-market tightness and wage setting. The results below therefore do not claim that a partial-equilibrium formula automatically extends to the macroeconomic target.

\section*{An implementable heterogeneous insurance benchmark}
Let an unemployed type have assets $a_\theta>0$, welfare weight $w_\theta\geq0$, and an exogenous unemployment probability $r_\theta\in[0,1]$. The insurer observes only $X=x$, a finite partition of types. Choose $0\leq b_x\leq\bar b$. Treat the shadow public-fund value $\lambda>0$ as fixed, corresponding to a specified outside financing instrument or a Lagrange multiplier on expected benefits. Employment utility is unaffected by $b$ in this benchmark. The benefit-dependent objective is
\[
 J(b)=\E\left[r_\theta\{w_\theta u(a_\theta+b_X;\theta)-\lambda b_X\}\right].
\]
Suppose $u$ is differentiable, increasing and strictly concave in consumption, and the weighted objective is strictly concave in each populated group. Then each $b_x$ is the unique maximizer of a one-dimensional concave function. An interior choice satisfies
\[
 \E[w_\theta u_c(a_\theta+b_x;\theta)\mid X=x,U]=\lambda, \tag{1}
\]
where the conditional type distribution among the unemployed is reweighted by $r_\theta$. At $b_x=0$ the left-hand side of (1) must be at most $\lambda$; at $b_x=\bar b$ it must be at least $\lambda$. These are sufficient as well as necessary by concavity.

For a homogeneous group with logarithmic utility and weight $w>0$, the rule is
\[
 b_x=\min\{\bar b,\max\{0,w/\lambda-a\}\}.
\]
For heterogeneous assets within a group the average of $w/(a+b)$ appears in (1); substituting average assets into the homogeneous formula is generally wrong. Observationally identical groups with different asset distributions can have different optimal benefits even if their unemployment-duration elasticities agree.

\section*{The value and feasibility of additional information}
If $X'$ refines $X$, every $X$-measurable benefit rule is also $X'$-measurable. Therefore optimized value weakly rises, under the same expected budget or public-fund criterion. In the separable fixed-$\lambda$ benchmark this can also be seen from
\[
 \E\max_b\E[j_\theta(b)\mid X']\geq
 \E\max_b\E[j_\theta(b)\mid X],
 \quad j_\theta(b)=r_\theta\{w_\theta u(a_\theta+b;\theta)-\lambda b\}.
\]
A strict gain requires groups merged by $X$ to prefer genuinely different feasible choices and have positive mass. If they share the same maximizer, extra information adds no value.

Duration, observed earnings history and aggregate state can be compared by treating each as an additional refinement, but this set-inclusion argument alone gives no additive decomposition: the gain from duration can depend on whether history is already observed. Conditioning on a private type would give only an unattainable upper benchmark unless truthful disclosure and its incentives are added. Legal admissibility is likewise a constraint, not something the value calculation can waive.

\section*{An exact local liquidity/search calculation}
Now consider a distinct one-period search model with effort $e$, job-finding probability $p(e)$ and disutility $k(e)$. A worker takes a payroll tax $t$ and wage $w$ as fixed and chooses
\[
 U(e;a,b,t)=(1-p(e))u(a+b)+p(e)u(a+w-t)-k(e).
\]
Assume an interior locally unique optimum, $p'>0$, and second derivative
$F_e=p''(u_E-u_U)-k''<0$, where $F=p'(u_E-u_U)-k'=0$ is the first-order condition. The implicit-function theorem gives
\[
 e_b=\frac{p'u_U'}{F_e},\qquad
 e_a=\frac{p'(u_U'-u_E')}{F_e},\qquad
 e_b-e_a=\frac{p'u_E'}{F_e}. \tag{2}
\]
The benefit increase raises only unemployment consumption; an asset increase raises consumption in both employment states. Their difference isolates a component that does not arise from the common wealth increase. With $u_U'>u_E'$ all three displayed responses have their stated negative signs. These formulas are a local model calculation, not a reproduction of a sufficient-statistics formula from the Chetty paper or its erratum.

Observing only $p'e_b$ does not separately identify $e_a$, the utility derivatives or effort curvature. A liquidity intervention affecting assets but not incentives, together with its exclusion restrictions, is needed to compare the responses in (2). Moreover, equilibrium wages, tax finance, job availability and acceptance behavior add channels beyond this worker's fixed-$t$ experiment.

\section*{Why balancing the payroll budget changes the derivative}
For a homogeneous population, suppose taxes are paid only in employment and balance benefits, so $t=(1-p)b/p$ with $0<p<1$. Along a differentiable symmetric equilibrium, let $p'$ now denote the total derivative of employment with respect to benefits, including induced taxes. At an interior private effort optimum, the worker envelope calculation gives
\[
 \frac{dW}{db}=(1-p)u_U'-pu_E't'
 =(1-p)(u_U'-u_E')+\frac bp u_E'p'. \tag{3}
\]
The second equality differentiates the balanced-budget identity, giving
$t'=(1-p)/p-bp'/p^2$. Thus even this simplified closure needs the equilibrium employment derivative rather than the partial derivative holding taxes fixed. At $p\downarrow0$ the financing rule itself becomes singular, so (3) cannot justify a uniform policy near zero employment. With aggregate matching externalities, the private effort envelope does not remove all welfare effects of changes in market tightness.

\section*{Robust policy rankings over incomplete models}
For a finite collection of admissible dynamic models $m=1,\ldots,M$ and a finite legally feasible schedule menu $b^1,\ldots,b^K$, compute each model's equilibrium welfare and net budget using its own stated closure. A robust deterministic choice maximizes
$\min_m W_m(b^k)$ over schedules satisfying every required model-specific budget. Pairwise dominance is certified only if $W_m(b^k)\geq W_m(b^h)$ for every $m$. A local average reform response does not identify the whole welfare table.

If public randomization over schedules before outcomes is explicitly allowed, probabilities $z_k$ give the linear program
\[
 \max_{z,t}\ t,\quad \sum_kz_kW_m(b^k)\geq t\ (m=1,\ldots,M),
 \quad z_k\geq0,\quad\sum_kz_k=1,
\]
with corresponding expected-budget inequalities. This is a different policy class from individual random benefit payments or a deterministic schedule; its gain should not be attributed to personalization. Model uncertainty is kept fixed over an entire dynamic path.

\section{Completion Estimate}
31\%

This is a post hoc, heuristic estimate for the full catalogue target.
The attempt derives feasible static heterogeneous insurance rules, information-value ordering, local liquidity-search responses, and a financing derivative plus robust finite-menu comparisons. It does not solve or identify legally implementable dynamic benefit schedules with private types, duration, aggregate matching, wage adjustment, and equilibrium fiscal constraints.

\section*{Status and source checked}
The analysis establishes an implementable static rule, an information-value ordering, local search decompositions and a robust finite-menu formulation. It does not solve the full heterogeneous dynamic equilibrium or identify its primitives from existing reforms. The catalogue target remains unresolved.

Camille Landais, Pascal Michaillat and Emmanuel Saez (2018), \emph{A Macroeconomic Approach to Optimal Unemployment Insurance: Theory}, American Economic Journal: Economic Policy 10(2), 152--181. Publisher abstract inspected for the matching/tightness correction and its model dependence. \url{https://www.aeaweb.org/articles?id=10.1257/pol.20150088}.

\end{document}
